\subsection{The doily code as a point-derived bent biplane}
\label{subsec:bent-biplane-klein}

The quadratic-evaluation code of the doily admits a useful one-coordinate completion.  On
$V=\mathbb F_2^4$ take the standard symplectic form $B$ and
$q_0(x)=x_0x_1+x_2x_3$.  The previously certified punctured code is
\[
 C=\{(B(a,x)+tq_0(x))_{x\ne0}:a\in V,\ t\in\mathbb F_2\}
   =[15,5,6]_2.
\]
Restoring the zero coordinate and the affine constant gives
\[
 D=\langle RM(1,4),q_0\rangle=[16,6,6]_2,
 \qquad
 W_D=1+16y^6+30y^8+16y^{10}+y^{16}.
\]
Both $C$ and $D$ are self-orthogonal.  Their generalized Hamming weights are
$(6,10,12,14,15)$ and $(6,10,12,14,15,16)$ respectively, and both have
covering radius six.

The sixteen minimum supports of $D$ form a symmetric $2$-$(16,6,2)$ biplane:
every two blocks meet in two points.  Fixing the restored zero coordinate recovers the
entire doily Veldkamp-point taxonomy.  The ten blocks avoiding the distinguished point
shorten to the ten weight-six grid-complement words of $C$; the six blocks through it,
after deleting that point, are the six ovoids; and the thirty weight-eight words pair by
complement, with the fifteen representatives vanishing at the distinguished coordinate
shortening to the fifteen perp words.  Thus one point of the biplane derives
\[
 10\text{ grids}+15\text{ perps}+6\text{ ovoids}=31.
\]
The full biplane automorphism group is
\[
 \operatorname{Aut}(\mathcal B)=2^4\!:\operatorname{Sp}(4,2)
 \cong2^4\!:\!S_6,
 \qquad |\operatorname{Aut}(\mathcal B)|=11520.
\]

There is simultaneously a Klein-geometric description of the earlier $15+20$ line split.
The $35$ lines of $PG(3,2)$ are the $35$ binary points of the Klein quadric
\[
 p_{01}p_{23}+p_{02}p_{13}+p_{03}p_{12}=0\subset PG(5,2).
\]
For the standard symplectic form, isotropy is exactly the hyperplane condition
\[
 p_{01}+p_{23}=0,
\]
so its section consists of the $15$ doily lines and the remaining $20$ Klein points are the
nonisotropic lines.  Symplectic polarity pairs those $20$ lines into ten skew pairs; their
six-point unions are exactly the ten grid-complement supports above.  If $M$ is their
$10\times15$ incidence matrix, then
\[
 MM^{\mathsf T}=4I_{10}+2J_{10},
 \qquad
 \operatorname{Spec}(MM^{\mathsf T})=\{24^1,4^9\}.
\]
This is a finite coding/design/projective-geometric theorem only; no physical identification
is inferred from it.
